\documentclass[10pt,a4paper]{amsart}
\usepackage[margin=19mm]{geometry}
\usepackage{amsmath,amssymb,mathtools}
\usepackage[scr=rsfs,cal=euler]{mathalfa}
\usepackage{xfrac}
\usepackage[unicode,hidelinks,bookmarks=false]{hyperref}
\newcommand{\ie}{i.e.\ }
\newcommand{\cf}{cf.\ }
\newcommand{\resp}{respectively\ }
\newcommand{\Subsection}{\subsection}
\providecommand{\nobreakdash}{\nobreak\hbox{-}}
\setlength{\emergencystretch}{2em}
\multlinegap0pt
\usepackage{enumerate}
\usepackage{tikz}
\usepackage{amssymb}
\usepackage{xcolor}
\newtheorem{lemma}{Lemma}
 \newcommand{\ba}{\begin{aligned}}
\newcommand{\be}{\begin{equation}}
 \newcommand{\ea}{\end{aligned}}
\newcommand{\ee}{\end{equation}}
  \newcommand{\f}{\frac}
\newcommand{\s}{\mathrm{div}}
\title{On energy equality of the 3D anisotropic Navier-Stokes equations}
\author{Yulin Ye, Wei Wei, Yanqing Wang}
\date{Source: arXiv:2609.12383v1 (CC BY 4.0)}
\begin{document}
\maketitle

\noindent\emph{Source and licence.} On energy equality of the 3D anisotropic Navier-Stokes equations. Version arXiv:2609.12383v1 (CC BY 4.0), distributed by arXiv under the Creative Commons Attribution 4.0 International licence, http://creativecommons.org/licenses/by/4.0/, as recorded in the arXiv licence metadata for this exact version and shown on the abstract page https://arxiv.org/abs/2609.12383v1. Full licence notice: \texttt{licenses/arxiv-2609.12383-CC-BY-4.0.txt}.\par\smallskip

\noindent\textit{Editorial scope.} Counted unit: the article's lemma pLions (source0.tex lines 273--279). The article's complete original proof of it is delivered with it (source0.tex lines 280--346, one \texttt{proof} environment of 67 lines). No mathematics is written, repaired or re-derived: the statement and the proof are the article's own lines, in the article's own notation, and no retained line is substituted or re-flowed.  No further source-local context is needed: every cross-reference of every retained line resolves inside the two delivered spans.  Nothing was omitted from any retained span, so the package carries no omission record.  No retained line contains a citation, so the wrapper carries no bibliography and no bibliography entry.  No retained line contains a figure environment, an image-inclusion command or drawing code, so no image is delivered and the package declares no asset dependency.  Editorial formatting only, in addition: an AMS \texttt{amsart} preamble that restates the article's own declarations and macros used by the retained lines, namely the environments \texttt{lemma} on separate counters, together with the standard packages those lines need.  The retained environments print in this document as Lemma 1 in order of appearance, whereas the article prints the same items with the numbering of its own full document.  The complete native source of the article as distributed in the arXiv e-print for this exact version is bundled unchanged as \texttt{sources/210111/source0.tex}.
\begin{lemma}
	\label{pLions}  Assume that the exponents satisfy $$1\leq p_1,q_1\leq \infty,~1\leq p, q, p_k, q_k< \infty,~k=2,3,4,$$ and the scaling relations $$\frac{1}{p}=\frac{1}{p_1}+\frac{1}{p_2}=\frac{1}{p_3}+\frac{1}{p_4}~ and ~ \frac{1}{q}=\frac{1}{q_1}+\frac{1}{q_2}=\frac{1}{q_3}+\frac{1}{q_4}.$$ Suppose that $$\partial_{i}f\in L^{p_2}(0,T;L^{q_2}(\mathbb{R}^3)), ~\text{and}~f\in L^{p_4}(0,T;L^{q_4}(\mathbb{R}^3)),$$
	and the divergence-free vector field $u$ fulfills
	 $$  u\in L^{p_1}(0,T;L^{q_1}(\mathbb{R}^3)),~\text{and}~\partial_{i} u\in L^{p_3}(0,T;L^{q_3}(\mathbb{R}^3)),$$
	 for all $i=1, 2,\cdots,n-1.$
Then,  $$  \s[(uf)^\varepsilon  -   u  f^\varepsilon   ]\to 0\quad\text{ strongly~in } {L^{p}(0,T;L^{q}(\mathbb{R}^3))},~\text{as}~\varepsilon\to 0.$$
\end{lemma}

\begin{proof}
First, by direct computation, we deduce that
	\be\label{2.5}\ba \s[(uf)^\varepsilon  -   u  f^\varepsilon   ]  =&  \sum_{i=1}^{  n-1}\partial_{i}[(u_{i}f)^\varepsilon  -   u_{i}  f^\varepsilon   ] +\partial_{n}[(u_{n}f)^\varepsilon  -   u_{n}  f^\varepsilon   ]\\ = & \sum_{i=1}^{  n-1}[\partial_{i}(u_{i}f)^\varepsilon  -   u_{i}  \partial_{i}f^\varepsilon   ] +[\partial_{n}(u_{n}f)^\varepsilon  -   u_{n}\partial_{n}  f^\varepsilon   ]\\ = & \sum_{i=1}^{  n-1}[(\partial_{i}u_{i}f)^\varepsilon + (u_{i}\partial_{i}f)^\varepsilon-   u_{i}  \partial_{i}f^\varepsilon   ] +[\partial_{n}(u_{n}f)^\varepsilon  -   u_{n}\partial_{n}  f^\varepsilon   ]\\
	=&R_1^\varepsilon +R_2^\varepsilon,\ea\ee
	where the divergence-free condition $\s u=0$ has been used.
	
	For the first term $R_1^\varepsilon$, it can be reformulated as
	\be\label{2.6}
	\ba
	R_1^\varepsilon&=\sum_{i=1}^{  n-1}[(\partial_{i}u_{i}f)^\varepsilon + (u_{i}\partial_{i}f)^\varepsilon-   u_{i}  \partial_{i}f^\varepsilon   ]\\
	&=\sum_{i=1}^{  n-1}(\partial_i u_i f)^\varepsilon+ \sum_{i=1}^{  n-1}[(u_i \partial_i f)^\varepsilon -(u_i \partial_i f)+(u_i \partial_i f)-u_i \partial_i f^\varepsilon].
		\ea\ee
Regarding the second term on the right-hand side of \eqref{2.6}, it follows from the H\"older inequality that
\be\label{2.7}\ba
\|u_i \partial_i f\|_{L^p(0,T;L^q(\mathbb{R}^3))}\leq \|u_i\|_{L^{p_1}(0,T;L^{q_1}(\mathbb{R}^3))}\|\partial_i f\|_{L^{p_2}(0,T;L^{q_2}(\mathbb{R}^3))},
\ea\ee
which together with the standard properties of the mollification yields
\be\label{2.8}\ba
\|(u_i \partial_i f)^\varepsilon - u_i \partial_i f\|_{L^p(0,T;L^q(\mathbb{R}^3))}\rightarrow 0,~~\text{as}~~\varepsilon \to 0,
 \ea\ee	
and
\be\label{2.9}\ba
&\|(u_i \partial_i f)-u_i \partial_i f^\varepsilon\|_{L^p(0,T;L^q(\mathbb{R}^3))}\\
\leq& \|u_i\|_{L^{p_1}(0,T;L^{q_1}(\mathbb{R}^3))}\|\partial_i f^\varepsilon -\partial_i f\|_{L^{p_2}(0,T;L^{q_2}(\mathbb{R}^3))}\to 0, ~~\text{as}~~\varepsilon \to 0,
\ea\ee
provided that $1\leq p,q, p_1, q_1\leq \infty, 1\leq p_2,q_2<\infty$ satisfying $\f1p=\frac{1}{p_1}+\frac{1}{p_2}$ and $\frac{1}{q}=\frac{1}{q_1}+\frac{1}{q_2}$.

Consequently, in combination with \eqref{2.6}, \eqref{2.8} and \eqref{2.9}, we conclude that
\be\label{2.10}
R_1^\varepsilon \rightarrow \sum_{i=1}^{n-1} \partial_i u_i f~~\text{strongly~in}~L^p(0,T;L^q(\mathbb{R}^3)),~\text{as}~\varepsilon \to 0.
\ee
For the second term $R_2^\varepsilon$, in light of the mean value theorem, we get
\be\label{2.11}\ba
R_2^\varepsilon &=\partial_{n}(u_{n}f)^\varepsilon  -   u_{n}\partial_{n}  f^\varepsilon\\
&=\int_{\mathbb{R}^3}\partial_{y_n} \eta_\varepsilon (y)\left(u_n(x-y)-u_n(x)\right)f(x-y) dy\\
&=\f{1}{\varepsilon}\int_{\mathbb{R}^3}\partial_{z_n}\eta(z)\left(u_n(x-\varepsilon z)-u_n(x)\right) f(x-\varepsilon z) dz\\
&=-\f{1}{\varepsilon}\int_{\mathbb{R}^3}\partial_{z_n}\eta(z)\int_0^1 \varepsilon z \cdot \nabla u_{n}(x-\theta \varepsilon z) d\theta  f(x-\varepsilon z) dz\\
&=-\sum_{i=1}^n \int_{\mathbb{R}^3}\partial_{z_n}\eta(z) z_i\int_0^1 \partial_i u_{n}(x-\theta \varepsilon z) d\theta  f(x-\varepsilon z) dz.
\ea\ee
On the one hand, by virtue of the Minkowski inequality and the H\"older inequality, we deduce that
\be\label{2.12}\ba
&\|\partial_{n}(u_{n}f)^\varepsilon  -   u_{n}\partial_{n}  f^\varepsilon\|_{L^p(0,T;L^q(\mathbb{R}^3)}\\
\leq & \sum_{i=1}^n \int_{\mathbb{R}^3}\int_0^1 |\partial_{z_n}\eta (z) z_i |\|\partial_i u_n(x-\theta \varepsilon z)\|_{L^{p_3}(0,T;L^{q_3}(\mathbb{R}^3))}\|f(x-\varepsilon z)\|_{L^{p_4}(0,T;L^{q_4}(\mathbb{R}^3))} d\theta dz\\
\leq& C\|\partial_i u_n\|_{L^{p_3}(0,T;L^{q_3}(\mathbb{R}^3))}\|f\|_{L^{p_4}(0,T;L^{q_4}(\mathbb{R}^3))},
\ea\ee
where $\f{1}{p}=\f{1}{p_3}+\f{1}{p_4}$ and $\f{1}{q}=\f{1}{q_3}+\f{1}{q_4}$.\\
On the other hand, the  continuity of the translations in Lebergue space gives
\be\label{2.13}\ba
&\partial_i u_n(x-\theta \varepsilon z)\rightarrow \partial_i u_n(x)~&\text{strongly~in~} L^{p_3}(0,T;L^{q_3}(\mathbb{R}^3)), ~\text{as} ~\varepsilon \to 0,\\
&f(x-\theta \varepsilon z)\rightarrow f(x)~&\text{strongly~in~} L^{p_4}(0,T;L^{q_4}(\mathbb{R}^3)), ~\text{as} ~\varepsilon \to 0,
\ea\ee
provided that $1\leq p_3,q_3,p_4,q_4<\infty$.

Then combining \eqref{2.11},\eqref{2.12} and \eqref{2.13}, we conclude from the dominated convergence theorem that
\be\label{2.14}\ba
R_2^\varepsilon &=\partial_{n}(u_{n}f)^\varepsilon  -   u_{n}\partial_{n}  f^\varepsilon\\
&\to -\sum_{i=1}^n \left(\int_{\mathbb{R}^3} \partial_{z_n} \eta(z) z_i  dz\right) \partial_i u_n(x) f(x)~&\text{strongly~in~} L^{p}(0,T;L^{q}(\mathbb{R}^3)), ~\text{as} ~\varepsilon \to 0\\
&=\partial_n u_n(x) f(x).
\ea\ee
Here in the last equality, we have used the integration by parts and the facts that $\sum_{i=1}^n \partial_{z_n} z_i=1$ and $\int_{\mathbb{R}^3} \eta(z) dz=1$.

Thus, in combination with \eqref{2.10} and \eqref{2.14}, and using the divergence-free condition, we obtain
\be\label{2.15}\ba
R_1^\varepsilon+R_2^\varepsilon &\to (\sum_{i=1}^{n-1} \partial_i u_i+\partial_n u_n) f=0, ~&\text{strongly~in~} L^{p}(0,T;L^{q}(\mathbb{R}^3)), ~\text{as} ~\varepsilon \to 0.
\ea\ee
Thus, the proof of this lemma is completed.
\end{proof}

\end{document}
